\documentclass[11pt,letterpaper]{article}
\usepackage[T1]{fontenc}
\usepackage[utf8]{inputenc}
\usepackage{lmodern}
\usepackage[margin=1in]{geometry}
\usepackage{amsmath,amssymb,amsthm,mathtools}
\usepackage{microtype,booktabs,array,longtable}
\usepackage{enumitem,listings,fancyhdr}
\usepackage[hidelinks]{hyperref}
\hypersetup{
  pdftitle={Fixed-multiplicity shuffles and balanced Smirnov words},
  pdfsubject={The exp(-(k-1)) law, Poisson adjacency statistics, and OEIS sequences}
}
\setlength{\parindent}{0pt}
\setlength{\parskip}{0.46em}
\setlength{\headheight}{14pt}
\pagestyle{fancy}
\fancyhf{}
\fancyhead[L]{\small Fixed-multiplicity shuffles and balanced Smirnov words}
\fancyhead[R]{\small\thepage}
\renewcommand{\headrulewidth}{0.3pt}

\newtheorem{theorem}{Theorem}[section]
\newtheorem{proposition}[theorem]{Proposition}
\newtheorem{lemma}[theorem]{Lemma}
\newtheorem{corollary}[theorem]{Corollary}
\newtheorem{conjecture}[theorem]{Conjecture}
\theoremstyle{definition}
\newtheorem{definition}[theorem]{Definition}
\theoremstyle{remark}
\newtheorem{remark}[theorem]{Remark}

\newcommand{\NN}{\mathbb N}
\newcommand{\CC}{\mathbb C}
\newcommand{\EE}{\mathbb E}
\newcommand{\PP}{\mathbb P}
\newcommand{\ee}{\mathrm e}
\newcommand{\fall}[2]{(#1)_{\underline{#2}}}
\newcommand{\coef}[2]{[#1]#2}
\newcommand{\Pois}{\operatorname{Poisson}}
\newcommand{\Var}{\operatorname{Var}}
\newcommand{\cO}{\mathcal O}
\newcommand{\W}{\operatorname{W}}
\DeclareMathOperator{\GammaDist}{Gamma}

\lstset{
  language=Python,
  basicstyle=\small\ttfamily,
  columns=fullflexible,
  keepspaces=true,
  showstringspaces=false,
  breaklines=true,
  frame=single,
  framerule=0.3pt,
  xleftmargin=0.8em,
  xrightmargin=0.8em,
  aboveskip=0.9em,
  belowskip=0.9em
}

\title{Fixed-multiplicity shuffles and balanced Smirnov words\\[0.45em]
\large The $\ee^{-(k-1)}$ law, Poisson adjacency statistics,\\
and explicit all-orders corrections}
\author{A research note on OEIS A007060, A193624, A330266, and related sequences}
\date{1 October 2026}

\begin{document}
\maketitle
\thispagestyle{empty}

\begin{abstract}
Let $B_{n,k}$ be the number of permutations of $kn$ distinct cards, partitioned
into $n$ ranks of size $k$, in which neighboring cards never have the same
rank.  The OEIS entry A330266~\cite{oeis330266} records the fixed-$k$
conjecture (S.~Kirgizov, 29 September 2023)
\[
  \frac{B_{n,k}}{(kn)!}\longrightarrow \ee^{-(k-1)}.
\]
The case $k=4$ was subsequently deduced on that entry from an approximation
to a related sequence (D.~Radcliffe, 9 September 2025), but the general
statement remained posed there.  We
give a direct proof for every fixed $k\ge2$ and prove considerably more.
If $X_{n,k}$ is the number of equal-rank adjacencies in a uniformly random
shuffle, then
\[
  X_{n,k}\xrightarrow{d}\Pois(k-1).
\]
We derive an exact probability-generating integral whose kernel is an
associated Laguerre polynomial, an exact inclusion--exclusion formula, a
gamma-expectation representation, and a constructive expansion to arbitrary
order.  In particular, uniformly for fixed $k$,
\[
\begin{aligned}
\frac{B_{n,k}}{(kn)!}=\ee^{-(k-1)}\bigg[&
1-\frac{(k-1)^2}{2kn}
+\frac{(k-1)^2(3k^2-14k+7)}{24k^2n^2}\\
&-\frac{(k-1)^4(k^2-10k+17)}{48k^3n^3}
+O_k(n^{-4})\bigg].
\end{aligned}
\]
The first four factorial cumulants are obtained in exact rational form; the
mean is exactly $k-1$ and the variance is
$(k-1)-(k-1)^2/(kn-1)$.  For A330266 this gives
\[
\frac{B_{n,4}}{(4n)!}=\ee^{-3}\left(1-\frac{9}{8n}
-\frac{3}{128n^2}+\frac{189}{1024n^3}+O(n^{-4})\right).
\]
We also treat the corresponding multiset sequences, local probabilities,
Hamiltonian paths in complete multipartite graphs, ratio asymptotics, and a
Lambert-$W$ inverse transseries.  Exact reproducibility code accompanies the
article.
\end{abstract}

\vfill
\small
\textbf{Scope and priority.}
The sequence definitions and historical comments below are taken from the
OEIS pages cited in the bibliography.  A330266~\cite{oeis330266} now states that its particular
$k=4$ leading equivalent follows from an approximation to A190830.  The point
of the present note is the uniform fixed-$k$ theorem, the Poisson law, exact
factorial cumulants, and systematic correction terms.  No claim of a complete
literature search or historical priority is made; the results should be
independently checked before submission or attribution.
\normalsize

\newpage
{\small\textbf{Editorial note (1 October 2026, batch 73O2).}
Manuscript~57 of batch~73 of ProveIt's incoming-reports intake, placed in
commit \texttt{6e193dd4f}; AI-assisted, unrefereed, not formalized.  The
write added the label prefix \texttt{bsw:}, citation commands, the names of
the OEIS contributors, notes marked ``[Added 1 October 2026, batch 73O2]'',
Remark~\ref{bsw:rem:tail-sketch}, one bibliography entry and
Appendix~\ref{bsw:app:provenance}.  The analytic tail step of
Theorem~\ref{bsw:thm:all-orders} is a sketch; the fixed-$k$ limit law does
not depend on it.\par}
\vspace{-0.5em}
\setcounter{tocdepth}{1}
\tableofcontents
\newpage

\section{The OEIS problem and the family of sequences}

Fix integers $n\ge0$ and $k\ge2$.  There are $n$ ranks, each containing $k$
distinct cards.  A shuffle is a permutation of all $N=kn$ cards.  Let
\begin{equation}\label{bsw:eq:Xdef}
 X_{n,k}=\#\{1\le j<N:\text{the cards at positions $j$ and $j+1$
have the same rank}\}.
\end{equation}
Write
\begin{equation}\label{bsw:eq:defB}
 B_{n,k}=N!\,\PP(X_{n,k}=0),
 \qquad p_{n,k}=\frac{B_{n,k}}{N!}.
\end{equation}
If cards within one rank are made indistinguishable, let $S_{n,k}$ denote the
number of admissible multiset words.  Every rank word lifts to exactly
$(k!)^n$ labeled permutations, so
\begin{equation}\label{bsw:eq:labeled-multiset}
 B_{n,k}=(k!)^n S_{n,k},
 \qquad
 p_{n,k}=\frac{S_{n,k}}{(kn)!/(k!)^n}.
\end{equation}
Thus the same probability governs the labeled and multiset versions.

The most visible OEIS rows are as follows.
\begin{center}
\begin{tabular}{ccll}
\toprule
$k$ & labeled count $B_{n,k}$ & interpretation & multiset count $S_{n,k}$\\
\midrule
2 & A007060~\cite{oeis007060} & couples with no adjacent spouses & A114938~\cite{oeis114938}\\
3 & A193624~\cite{oeis193624} & triples with no adjacent siblings & A193638~\cite{oeis193638}\\
4 & A330266~\cite{oeis330266} & four cards in each rank & A321633~\cite{oeis321633}\\
\bottomrule
\end{tabular}
\end{center}
A330266~\cite{oeis330266} records the exact integral
\[
 B_{n,4}=\int_0^\infty
 (t^4-12t^3+36t^2-24t)^n\ee^{-t}\,dt
\]
and the conjectured fixed-multiplicity law, posed there by S.~Kirgizov on
29 September 2023,
\begin{equation}\label{bsw:eq:oeis-conj}
 p_{n,k}\longrightarrow \ee^{-(k-1)}.
\end{equation}
The entry also records a September 2025 deduction of the $k=4$ equivalent
from a separate approximation: D.~Radcliffe (9 September 2025) combined
$\textnormal{A330266}(n)=24^n\,n!\,\textnormal{A190830}(n)$ with
V.~Kot\v{e}\v{s}ovec's approximation of A190830 and Stirling's formula.  Our derivation below is direct, applies to
every fixed $k$, and identifies the full limiting distribution behind
\eqref{bsw:eq:oeis-conj}.

\subsection{Why a Poisson law is plausible but not automatic}

For each of the $N-1$ adjacent position pairs, the probability of equal ranks
is $(k-1)/(N-1)$.  Hence
\begin{equation}\label{bsw:eq:mean-quick}
 \EE X_{n,k}=(N-1)\frac{k-1}{N-1}=k-1.
\end{equation}
If the rare adjacency indicators were independent, one would expect a
Poisson limit of mean $k-1$ and hence a zero probability $\ee^{-(k-1)}$.
They are not independent: overlapping position pairs and the fixed rank
multiplicities produce both local and global dependence.  The proof therefore
requires more than the heuristic ``rare events with bounded mean.''  The
exact generating function will make the cancellation transparent.

\subsection{Hamiltonian-path interpretation}

Replace the deck by the complete $n$-partite graph with $n$ parts of size
$k$.  A vertex ordering is an admissible shuffle precisely when consecutive
vertices are joined.  Consequently $B_{n,k}$ is the number of \emph{oriented}
Hamiltonian paths, counted as ordered vertex lists, in
$K_{k,k,\ldots,k}$.  Unoriented paths are counted by $B_{n,k}/2$ when
$kn>1$.  This connection places the problem among balanced Smirnov words and
complete multipartite Hamiltonian enumeration~\cite{eriksson-martin,mehiri}.

\section{A weighted run transform}

It is useful to enumerate all shuffles while recording the number of bad
adjacencies.  Define
\begin{equation}\label{bsw:eq:weightedA}
 A_{n,k}(u)=\sum_{\pi\in\mathfrak S_{kn}}u^{X_{n,k}(\pi)},
 \qquad
 G_{n,k}(u)=\frac{A_{n,k}(u)}{(kn)!}=\EE u^{X_{n,k}}.
\end{equation}
Thus $A_{n,k}(0)=B_{n,k}$ and $G_{n,k}(0)=p_{n,k}$.

\begin{theorem}[Run-transform generating function]\label{bsw:thm:run-transform}
For indeterminates $x_1,\ldots,x_n$,
\begin{equation}\label{bsw:eq:runGF}
 \sum_{w}u^{X(w)}x_1^{|w|_1}\cdots x_n^{|w|_n}
 =\frac{1}{1-\displaystyle\sum_{i=1}^n
             \frac{x_i}{1+(1-u)x_i}},
\end{equation}
where the sum is over all finite words on $\{1,\ldots,n\}$, including the
empty word.  Therefore
\begin{equation}\label{bsw:eq:coeffA}
 A_{n,k}(u)=(k!)^n\coef{x_1^k\cdots x_n^k}
 \frac{1}{1-\displaystyle\sum_{i=1}^n
             \frac{x_i}{1+(1-u)x_i}}.
\end{equation}
\end{theorem}

\begin{proof}
Compress every maximal run of equal letters to one letter.  The compressed
word is a Smirnov word, meaning that consecutive letters are distinct.  A run
of letter $i$ and length $r\ge1$ contributes
$u^{r-1}x_i^r$, whose generating series is
\[
 y_i=\sum_{r\ge1}u^{r-1}x_i^r=\frac{x_i}{1-ux_i}.
\]
The classical multivariate generating function of Smirnov
words~\cite{stanley,flajolet-sedgewick} is
\[
 \frac{1}{1-\sum_i y_i/(1+y_i)}.
\]
Indeed, it follows from the ordinary word series
$1/(1-\sum_i z_i)$ by replacing each nonempty run by one marked letter, or
by the standard run-decomposition bijection.  Substitution gives
\[
 \frac{y_i}{1+y_i}=\frac{x_i}{1+(1-u)x_i},
\]
which proves \eqref{bsw:eq:runGF}.  Extracting $k$ occurrences of each rank gives
multiset words; labeling the $k$ copies independently in each rank contributes
$(k!)^n$ and proves \eqref{bsw:eq:coeffA}.
\end{proof}

\section{The exact Laguerre integral}

Set $v=u-1$ and define
\begin{equation}\label{bsw:eq:Pku}
 P_{k,u}(t)=k!\coef{x^k}
 \exp\!\left(\frac{tx}{1-vx}\right).
\end{equation}

\begin{theorem}[Exact integral and polynomial kernel]\label{bsw:thm:integral}
For every $n\ge0$, $k\ge2$, and complex $u$,
\begin{equation}\label{bsw:eq:exact-integral}
 \boxed{\displaystyle
 A_{n,k}(u)=\int_0^\infty \ee^{-t}P_{k,u}(t)^n\,dt.}
\end{equation}
The kernel has the explicit forms
\begin{align}
 P_{k,u}(t)
 &=\sum_{j=1}^k \binom{k-1}{j-1}\frac{k!}{j!}
       t^j(u-1)^{k-j},\label{bsw:eq:Pexplicit}\\
 &=k!(u-1)^k L_k^{(-1)}\!\left(-\frac{t}{u-1}\right),
 \label{bsw:eq:Plaguerre}
\end{align}
where the value at $u=1$ is interpreted by continuity.
\end{theorem}

\begin{proof}
Use the coefficientwise Laplace identity
\[
 \frac1{1-s}=\int_0^\infty \ee^{-t}\ee^{ts}\,dt.
\]
In \eqref{bsw:eq:coeffA}, the exponential factors into $n$ identical univariate
series.  Extracting $x_i^k$ from each factor gives
\eqref{bsw:eq:exact-integral}.  Expanding
$\exp(tx/(1-vx))$ first by the exponential degree and then extracting $x^k$
gives \eqref{bsw:eq:Pexplicit}.  Finally, the standard
expansion~\cite{dlmf}
\[
 L_k^{(-1)}(z)=\sum_{j=1}^k
 \binom{k-1}{j-1}\frac{(-z)^j}{j!}
\]
gives \eqref{bsw:eq:Plaguerre}.
\end{proof}

For no equal adjacencies, $u=0$ and $v=-1$.  Write
\begin{equation}\label{bsw:eq:Pk}
 P_k(t):=P_{k,0}(t)=(-1)^k k!L_k^{(-1)}(t).
\end{equation}
The first kernels are
\begin{align*}
 P_2(t)&=t^2-2t,\\
 P_3(t)&=t^3-6t^2+6t,\\
 P_4(t)&=t^4-12t^3+36t^2-24t,\\
 P_5(t)&=t^5-20t^4+120t^3-240t^2+120t.
\end{align*}
The third line recovers the integral printed on A330266, while the theorem
explains it as one row of a uniform Laguerre family.

\section{Exact inclusion--exclusion and a compact reciprocal polynomial}

For $0\le m\le k-1$, put
\begin{equation}\label{bsw:eq:akm}
 a_{k,m}=\binom{k-1}{m}\frac{k!}{(k-m)!}
        =m!\binom{k}{m}\binom{k-1}{m},
\end{equation}
and define the degree-$(k-1)$ polynomial
\begin{equation}\label{bsw:eq:Rk}
 R_k(z)=\sum_{m=0}^{k-1}a_{k,m}z^m.
\end{equation}
The no-adjacency reciprocal polynomial is
\begin{equation}\label{bsw:eq:Ck}
 C_k(z)=z^kP_k(1/z)=R_k(-z).
\end{equation}

\begin{theorem}[Exact finite formula]\label{bsw:thm:exact-sum}
Let $N=kn$.  Then
\begin{align}
 B_{n,k}
 &=\sum_{m_1,\ldots,m_n=0}^{k-1}
 (-1)^M(N-M)!\prod_{i=1}^n a_{k,m_i},
 \qquad M=m_1+\cdots+m_n,\label{bsw:eq:multi-sum}\\
 &=\sum_{M=0}^{n(k-1)}(N-M)!\coef{z^M}C_k(z)^n.
 \label{bsw:eq:grouped-sum}
\end{align}
More generally, the factorial probability generating function
\begin{equation}\label{bsw:eq:Phi-def}
 \Phi_{n,k}(v):=\EE(1+v)^{X_{n,k}}=G_{n,k}(1+v)
\end{equation}
satisfies the exact identity
\begin{equation}\label{bsw:eq:Phi-exact}
 \boxed{\displaystyle
 \Phi_{n,k}(v)=
 \sum_{M=0}^{n(k-1)}
 \frac{\coef{z^M}R_k(z)^n}{\fall{N}{M}}v^M.}
\end{equation}
\end{theorem}

\begin{proof}
From \eqref{bsw:eq:Pexplicit},
\[
 z^kP_{k,u}(1/z)=\sum_{m=0}^{k-1}a_{k,m}v^mz^m=R_k(vz).
\]
Raise to the $n$th power and integrate termwise in
\eqref{bsw:eq:exact-integral}.  Since
$\int_0^\infty t^{N-M}\ee^{-t}\,dt=(N-M)!$, division by $N!$ yields
\eqref{bsw:eq:Phi-exact}.  Setting $v=-1$ gives
\eqref{bsw:eq:multi-sum} and \eqref{bsw:eq:grouped-sum}.
\end{proof}

Formula \eqref{bsw:eq:grouped-sum} is a rook-polynomial form of
inclusion--exclusion.  It uses only a degree-$(k-1)$ polynomial raised to the
$n$th power and is efficient for exact integer computation.  In particular,
for $k=4$ it reproduces
\[
 B_{13,4}=
3668033946384704437729512814619767610579526911188666362431432294400,
\]
the standard-deck value recorded on A330266.

\section{The Poisson adjacency theorem}

Recall that the $r$th falling factorial is
$\fall{x}{r}=x(x-1)\cdots(x-r+1)$.

\begin{theorem}[Factorial moments]\label{bsw:thm:factorial-moments}
For fixed $k\ge2$ and fixed $r\ge0$,
\begin{equation}\label{bsw:eq:fm-limit}
 \EE\fall{X_{n,k}}{r}=(k-1)^r+O_{k,r}(n^{-1}).
\end{equation}
Consequently,
\begin{equation}\label{bsw:eq:poisson-limit}
 \boxed{\displaystyle
 X_{n,k}\xrightarrow[n\to\infty]{d}\Pois(k-1).}
\end{equation}
\end{theorem}

\begin{proof}
The factorial moment is $r!$ times the coefficient of $v^r$ in
\eqref{bsw:eq:Phi-exact}.  Since every term $a_{k,m}z^m$ in $R_k(z)$ has the same
$z$-degree as the corresponding $v$-degree, the denominator for this
coefficient is exactly $\fall{N}{r}$.  The leading contribution to
$\coef{z^r}R_k(z)^n$ selects $a_{k,1}z$ from $r$ distinct factors and $1$
from the rest:
\[
 \coef{z^r}R_k(z)^n
 =\binom nr a_{k,1}^r+O_{k,r}(n^{r-1}).
\]
Here $a_{k,1}=k(k-1)$.  Therefore
\begin{align*}
 \EE\fall{X_{n,k}}r
 &=\frac{r!\binom nr[k(k-1)]^r+O_{k,r}(n^{r-1})}
         {\fall{kn}{r}}\\
 &=(k-1)^r+O_{k,r}(n^{-1}).
\end{align*}
These are the factorial moments of $\Pois(k-1)$.  Ordinary moments are finite
linear combinations of factorial moments, so every ordinary moment converges.
The Poisson distribution is moment-determinate; the method of moments gives
\eqref{bsw:eq:poisson-limit}.
\end{proof}

\begin{corollary}[The fixed-multiplicity OEIS law]\label{bsw:cor:main-limit}
For every fixed integer $k\ge2$,
\begin{equation}\label{bsw:eq:main-limit}
 \boxed{\displaystyle
 \frac{B_{n,k}}{(kn)!}=p_{n,k}\longrightarrow\ee^{-(k-1)}.}
\end{equation}
Equivalently,
\begin{equation}\label{bsw:eq:multiset-leading}
 S_{n,k}\sim\frac{(kn)!}{(k!)^n}\ee^{-(k-1)}.
\end{equation}
\end{corollary}

\begin{proof}
Because the variables are integer-valued,
\[
 \PP(X_{n,k}=0)=\PP(-\tfrac12<X_{n,k}<\tfrac12).
\]
The endpoints $\pm\tfrac12$ carry no mass under the limiting Poisson law, so
weak convergence implies convergence of this probability.  Hence
$\PP(X_{n,k}=0)\to\PP(\Pois(k-1)=0)=\ee^{-(k-1)}$.
Equation \eqref{bsw:eq:multiset-leading} follows from
\eqref{bsw:eq:labeled-multiset}.
\end{proof}

\begin{remark}
The proof is short because the correct object is not merely the zero-event
count.  The entire adjacency statistic has a stable Poisson limit.  The
constant $k-1$ is already forced by the exact mean \eqref{bsw:eq:mean-quick}.
\end{remark}

\section{A gamma expectation and an all-orders expansion engine}

The exact finite sum admits a useful probabilistic transform.  Let
$T_N$ have the gamma density
\begin{equation}\label{bsw:eq:gamma-density}
 f_N(t)=\frac{t^N\ee^{-t}}{N!},\qquad t>0,
 \qquad T_N\sim\GammaDist(N+1,1).
\end{equation}
Since
\begin{equation}\label{bsw:eq:negative-moment}
 \EE T_N^{-M}=\frac{(N-M)!}{N!}=\frac1{\fall NM}
 \qquad(0\le M\le N),
\end{equation}
we obtain the following exact formula.

\begin{proposition}[Gamma representation]\label{bsw:prop:gamma}
For $N=kn$,
\begin{equation}\label{bsw:eq:gamma-rep}
 \boxed{\displaystyle
 \Phi_{n,k}(v)=\EE\left[R_k\!\left(\frac v{T_N}\right)^n\right].}
\end{equation}
\end{proposition}

The gamma variable is concentrated at scale $\sqrt N$ around $N+1$.  Since
$R_k(w)=1+k(k-1)w+O_k(w^2)$, the exponent in
\eqref{bsw:eq:gamma-rep} tends to $(k-1)v$.  This already supplies another proof
of the Poisson limit.  More importantly, it supports a complete Laplace
expansion.

For a nonnegative integer $M$, define polynomials $H_r(M)$ by
\begin{equation}\label{bsw:eq:Hr-def}
 \prod_{j=0}^{M-1}\frac1{1-jx}=\sum_{r\ge0}H_r(M)x^r.
\end{equation}
Thus $H_r(M)=h_r(0,1,\ldots,M-1)$, a complete homogeneous symmetric
polynomial.  The first four are
\begin{align}
 H_0(M)&=1,\label{bsw:eq:H0}\\
 H_1(M)&=\frac{M(M-1)}2,\label{bsw:eq:H1}\\
 H_2(M)&=\frac{M(M-1)(3M^2+M-2)}{24},\label{bsw:eq:H2}\\
 H_3(M)&=\frac{M^2(M-1)^2(M+1)(M+2)}{48}.
 \label{bsw:eq:H3}
\end{align}
Let $D=v\,d/dv$.

\begin{theorem}[Constructive all-orders expansion]\label{bsw:thm:all-orders}
Fix $k\ge2$.  As $n\to\infty$, uniformly for $v$ in every compact subset of
$\CC$, $\Phi_{n,k}(v)$ has a Poincare expansion in integer powers of $1/n$.
A coefficientwise construction is
\begin{equation}\label{bsw:eq:operator-expansion}
 \boxed{\displaystyle
 \Phi_{n,k}(v)\sim
 \sum_{r\ge0}\frac1{N^r}
 H_r(D)\left[R_k\!\left(\frac vN\right)^n\right],
 \qquad N=kn.}
\end{equation}
After expanding the bracket in $1/N$, collecting total powers gives any
desired correction term.  In particular, every coefficient is an explicit
polynomial in $v$ and a rational function of $k$, multiplied by
$\ee^{(k-1)v}$.
\end{theorem}

\begin{proof}
Write
\[
 R_k(z)^n=\sum_{M=0}^{n(k-1)}r_{n,M}z^M.
\]
Then \eqref{bsw:eq:Phi-exact} is
\[
 \Phi_{n,k}(v)=\sum_M r_{n,M}\frac{v^M}{N^M}
 \frac{N^M}{\fall NM}.
\]
For fixed $M$,
\[
 \frac{N^M}{\fall NM}
 =\prod_{j=0}^{M-1}\left(1-\frac jN\right)^{-1}
 \sim\sum_{r\ge0}\frac{H_r(M)}{N^r}.
\]
Since $D$ acts on $v^M$ by multiplication by $M$, this is precisely the
formal operator rule \eqref{bsw:eq:operator-expansion}.

It remains to justify the asymptotic use when $M$ is not fixed.  The exact
gamma representation \eqref{bsw:eq:gamma-rep} does so.  Split the integral into
$|T_N-N|\le N^{2/3}$ and its complement.  On the upper tail, gamma Chernoff
bounds are enough.  On the lower tail one combines the polynomial bound
$|R_k(v/t)|\le C_{k,V}(1+t^{1-k})$ for $|v|\le V$ with Stirling's formula;
the exact negative moments show that terms whose total polynomial degree is
proportional to $N$ are in fact superexponentially small.  Thus the weighted
complement is smaller than every fixed power of $N^{-1}$.  On the central
interval, $T_N^{-1}$, $\log R_k(v/T_N)$, and the gamma density all possess
Taylor expansions with remainders uniform for $v$ in a prescribed compact
set.
After the change $T_N=N+\sqrt N\,y$, each coefficient is a polynomial in $y$
times a Gaussian majorant, so termwise integration is valid to arbitrary
fixed order.  Odd half-powers cancel in the integrated expansion, leaving
integer powers of $1/N$.  The coefficientwise operator computation and the
Laplace computation agree because both expand the same exact negative moments
\eqref{bsw:eq:negative-moment}.  This proves the Poincare expansion and the stated
uniformity.
\end{proof}

\begin{remark}[Status of the analytic step]\label{bsw:rem:tail-sketch}
[Added 1 October 2026, batch 73O2.]  The tail estimate is sketched.  The
coefficientwise operator rule is exact algebra, but the passage to an
asymptotic statement uniform for $v$ in compact sets is not proved in detail:
the gamma Chernoff bounds, the superexponential smallness of terms of total
degree proportional to $N$, and the cancellation of odd half-powers are
asserted, not derived.  A uniform factorial-moment bound of the type proved in
the companion path-forest report \texttt{a189281-path-forest-expansions} of
this collection (its Lemma~4.1, ``Uniform factorial-moment bound'', label
\texttt{spf:lem:moment-bound}, which bounds the $k$th factorial moment by
$(\theta\ee^2)^k/k!$ for its own model) would make the truncation rigorous; that
lemma is not proved here for the present model.

Accordingly, Theorem~\ref{bsw:thm:all-orders} and everything that uses it
rest on this sketched step: the scale $\kappa_r^{(F)}=O_{k,r}(n^{1-r})$ stated
after Theorem~\ref{bsw:thm:cumulants}, Theorem~\ref{bsw:thm:log-pgf}
(through ``the fifth and higher factorial cumulants begin at
$O_k(n^{-4})$''), Corollaries~\ref{bsw:cor:p-expansion}
and~\ref{bsw:cor:local}, the specializations
\eqref{bsw:eq:k2}--\eqref{bsw:eq:k5}, the count, ratio and inverse expansions
\eqref{bsw:eq:logB}, \eqref{bsw:eq:ratio} and \eqref{bsw:eq:inverse}, and the
correction terms in the proposed OEIS comments.  They do not depend on it:
Theorems~\ref{bsw:thm:run-transform}, \ref{bsw:thm:integral},
\ref{bsw:thm:exact-sum}, \ref{bsw:thm:factorial-moments}
and~\ref{bsw:thm:cumulants}, Proposition~\ref{bsw:prop:gamma}, and
Corollary~\ref{bsw:cor:main-limit}, the fixed-$k$ limit law that settles the
conjecture on A330266.  The intake's numerical check is consistent with the
displayed coefficients: for $k=2,\ldots,5$ the residual after the
$n^{-3}$ term of \eqref{bsw:eq:pfull}, multiplied by $\ee^{k-1}n^4$, is
nearly the same at $n=100$ and $n=400$, as an $O(n^{-4})$ remainder requires.
That is evidence, not a proof.
\end{remark}

\section{Exact factorial cumulants}

Define factorial cumulants $\kappa_r^{(F)}$ by
\begin{equation}\label{bsw:eq:fact-cumulants}
 \log\Phi_{n,k}(v)=\sum_{r\ge1}\kappa_r^{(F)}\frac{v^r}{r!}.
\end{equation}
They are particularly natural for a Poisson limit: a Poisson variable of mean
$\lambda$ has $\kappa_1^{(F)}=\lambda$ and all higher factorial cumulants zero.

\begin{theorem}[First four factorial cumulants]\label{bsw:thm:cumulants}
Let $N=kn$.  Whenever the displayed denominators are nonzero,
\begin{align}
 \kappa_1^{(F)}&=k-1,\label{bsw:eq:kappa1}\\
 \kappa_2^{(F)}&=-\frac{(k-1)^2}{N-1},\label{bsw:eq:kappa2}\\
 \kappa_3^{(F)}&=\frac{2(k-2)(k-1)^2}{(N-2)(N-1)},\label{bsw:eq:kappa3}\\
 \kappa_4^{(F)}&=-\frac{6(k-1)^2
 \{N(k^2-6k+7)+4k-6\}}
 {(N-3)(N-2)(N-1)^2}.
 \label{bsw:eq:kappa4}
\end{align}
In particular,
\begin{equation}\label{bsw:eq:mean-var}
 \boxed{\displaystyle
 \EE X_{n,k}=k-1,
 \qquad
 \Var X_{n,k}=(k-1)-\frac{(k-1)^2}{kn-1}.}
\end{equation}
\end{theorem}

\begin{proof}
Keep terms through $v^4$ in \eqref{bsw:eq:Phi-exact}.  Since
\[
 R_k(z)=1+a_{k,1}z+a_{k,2}z^2+a_{k,3}z^3+a_{k,4}z^4+O(z^5),
\]
ordinary multinomial expansion gives the first four factorial moments as
rational functions of $n$ and $N$.  Taking the formal logarithm and using
\eqref{bsw:eq:akm} simplifies to \eqref{bsw:eq:kappa1}--\eqref{bsw:eq:kappa4}.  This is
finite polynomial algebra; Appendix~\ref{bsw:app:cumulant-algebra} records a
reproducible symbolic check.  For any integer-valued random variable,
$\EE X=\kappa_1^{(F)}$ and
$\Var X=\kappa_1^{(F)}+\kappa_2^{(F)}$, which gives
\eqref{bsw:eq:mean-var}.
\end{proof}

All higher factorial cumulants tend to zero.  The all-orders theorem gives the
sharper scale $\kappa_r^{(F)}=O_{k,r}(n^{1-r})$ for fixed $r\ge2$, a
mod-Poisson pattern that controls local corrections.

\section{Explicit expansion through third order}

Put $\lambda=k-1$.  Expanding the exact cumulants and applying
Theorem~\ref{bsw:thm:all-orders} gives the following compact logarithmic form.

\begin{theorem}[Logarithmic factorial-PGF expansion]\label{bsw:thm:log-pgf}
For fixed $k\ge2$, uniformly for $v$ in compact subsets of $\CC$,
\begin{align}
 \log\Phi_{n,k}(v)
={}&\lambda v
-\frac{\lambda^2v^2}{2kn}\notag\\
&+\frac{\lambda^2v^2\{2(k-2)v-3\}}{6k^2n^2}\notag\\
&-\frac{\lambda^2v^2
 \{(k^2-6k+7)v^2-4(k-2)v+2\}}
 {4k^3n^3}
+O_{k,v}(n^{-4}).
\label{bsw:eq:log-pgf-exp}
\end{align}
\end{theorem}

\begin{proof}
Insert \eqref{bsw:eq:kappa2}--\eqref{bsw:eq:kappa4} into
\eqref{bsw:eq:fact-cumulants} and expand their denominators in $1/n$.  These terms
supply every contribution through $n^{-3}$ because the fifth and higher
factorial cumulants begin at $O_k(n^{-4})$, as follows from
Theorem~\ref{bsw:thm:all-orders}.  Collecting powers of $v$ gives
\eqref{bsw:eq:log-pgf-exp}.
\end{proof}

Setting $v=-1$ produces the no-adjacency probability.

\begin{corollary}[No-adjacency transseries]\label{bsw:cor:p-expansion}
For fixed $k\ge2$,
\begin{align}
 \log p_{n,k}
={}&-(k-1)
-\frac{(k-1)^2}{2kn}
-\frac{(k-1)^2(2k-1)}{6k^2n^2}
-\frac{(k-1)^4}{4k^3n^3}
+O_k(n^{-4}),
\label{bsw:eq:logp}\\[0.5em]
 p_{n,k}
={}&\ee^{-(k-1)}\left[
 1-\frac{(k-1)^2}{2kn}
 +\frac{(k-1)^2(3k^2-14k+7)}{24k^2n^2}\right.\notag\\
&\hspace{8em}\left.
 -\frac{(k-1)^4(k^2-10k+17)}{48k^3n^3}
 +O_k(n^{-4})\right].
\label{bsw:eq:pfull}
\end{align}
\end{corollary}

The logarithmic form is often numerically superior because exponentiating a
truncated log preserves positivity and partially resums products of lower
orders.

\subsection{A first local correction}

The expansion describes more than the mass at zero.

\begin{corollary}[Fixed local probabilities]\label{bsw:cor:local}
For each fixed $j\ge0$,
\begin{equation}\label{bsw:eq:local}
 \PP(X_{n,k}=j)
 =\ee^{-\lambda}\frac{\lambda^j}{j!}
 \left[1-\frac{(\lambda-j)^2-j}{2kn}+O_{k,j}(n^{-2})\right].
\end{equation}
\end{corollary}

\begin{proof}
Exponentiating the first correction in \eqref{bsw:eq:log-pgf-exp} gives
\[
 \Phi_{n,k}(v)=\ee^{\lambda v}
 \left(1-\frac{\lambda^2v^2}{2kn}+O(n^{-2})\right).
\]
Write $v=u-1$ and extract the coefficient of $u^j$.  Dividing the correction
by the Poisson mass gives
\[
 -\frac{\lambda^2}{2kn}
 \left(1-\frac{2j}{\lambda}+\frac{j(j-1)}{\lambda^2}\right),
\]
which is \eqref{bsw:eq:local}.
\end{proof}

\section{Specializations to the OEIS rows}

For $k=2,3,4,5$, \eqref{bsw:eq:pfull} becomes
\begin{align}
 p_{n,2}&=\ee^{-1}\left(1-\frac1{4n}-\frac3{32n^2}
 -\frac1{384n^3}+O(n^{-4})\right),\label{bsw:eq:k2}\\
 p_{n,3}&=\ee^{-2}\left(1-\frac2{3n}-\frac4{27n^2}
 +\frac4{81n^3}+O(n^{-4})\right),\label{bsw:eq:k3}\\
 p_{n,4}&=\ee^{-3}\left(1-\frac9{8n}-\frac3{128n^2}
 +\frac{189}{1024n^3}+O(n^{-4})\right),\label{bsw:eq:k4}\\
 p_{n,5}&=\ee^{-4}\left(1-\frac8{5n}+\frac8{25n^2}
 +\frac{128}{375n^3}+O(n^{-4})\right).
\label{bsw:eq:k5}
\end{align}
Thus \eqref{bsw:eq:k2} sharpens the known $\ee^{-1}$ law for
A007060~\cite{oeis007060}, \eqref{bsw:eq:k3} sharpens the $\ee^{-2}$ law
printed on A193624~\cite{oeis193624}, and \eqref{bsw:eq:k4} supplies explicit
corrections for A330266~\cite{oeis330266}.

\subsection{The standard 52-card deck}

For $n=13$ and $k=4$, exact enumeration gives
\begin{equation}\label{bsw:eq:deck-prob}
 p_{13,4}=
 0.04547628233109430448968591097164118590821040881636\ldots.
\end{equation}
The successive approximations are
\begin{center}
\begin{tabular}{lr}
\toprule
approximation & numerical value\\
\midrule
$\ee^{-3}$ & $0.04978706836786394298$\\
through $n^{-1}$ in \eqref{bsw:eq:k4} & $0.04547857206679879407$\\
through $n^{-2}$ in \eqref{bsw:eq:k4} & $0.04547166742529067684$\\
through $n^{-3}$ in \eqref{bsw:eq:k4} & $0.04547585004466578632$\\
exponential of \eqref{bsw:eq:logp} through $n^{-3}$
& $0.04547622509184257975$\\
exact & $0.04547628233109430449$\\
\bottomrule
\end{tabular}
\end{center}
The third-order logarithmic approximation has relative error about
$1.26\times10^{-6}$ even at the modest value $n=13$.

The exact mean number of equal-rank adjacencies is $3$, while
\[
 \Var X_{13,4}=3-\frac9{51}=\frac{48}{17}=2.8235294117\ldots.
\]
This quantifies the slight underdispersion relative to a Poisson$(3)$ law.

\subsection{Numerical convergence across several multiplicities}

The following table lists $\ee^{k-1}p_{n,k}$, which tends to $1$.  The final
column is the signed error after subtracting all three displayed correction
terms in \eqref{bsw:eq:pfull}.
\begin{center}
\begin{tabular}{rrrr}
\toprule
$k$ & $n$ & $\ee^{k-1}p_{n,k}$ & third-order residual\\
\midrule
2&10&0.97406694843814&$7.05\times10^{-6}$\\
2&50&0.99496248890932&$9.74\times10^{-9}$\\
3&10&0.93191894503056&$1.77\times10^{-5}$\\
3&50&0.98660782736509&$2.49\times10^{-8}$\\
4&10&0.88747575415541&$2.56\times10^{-5}$\\
4&50&0.97749213853337&$3.70\times10^{-8}$\\
5&10&0.84356657118181&$2.52\times10^{-5}$\\
5&50&0.96813076805226&$3.74\times10^{-8}$\\
6&10&0.80107140086354&$1.29\times10^{-5}$\\
6&50&0.95869552482316&$1.94\times10^{-8}$\\
\bottomrule
\end{tabular}
\end{center}
The small residuals are not used as proof; they are independent diagnostics
of the symbolic expansion.

\section{Asymptotics of the counts, ratios, and inverse functions}

\subsection{Logarithmic count transseries}

Let $N=kn$ and $\lambda=k-1$.  Combining Stirling's expansion with
\eqref{bsw:eq:logp} yields
\begin{align}
 \log B_{n,k}
={}&N(\log N-1)+\frac12\log(2\pi N)-\lambda
 +\frac{A_1}{N}+\frac{A_2}{N^2}+\frac{A_3}{N^3}
 +O_k(N^{-4}),
\label{bsw:eq:logB}
\end{align}
where
\begin{equation}\label{bsw:eq:A123}
 A_1=\frac1{12}-\frac{\lambda^2}{2},\qquad
 A_2=-\frac{\lambda^2(2k-1)}6,\qquad
 A_3=-\frac1{360}-\frac{\lambda^4}{4}.
\end{equation}
For multiset words,
\begin{equation}\label{bsw:eq:logS}
 \log S_{n,k}=\log B_{n,k}-n\log(k!).
\end{equation}
These formulas separate the universal factorial growth from the finite
adjacency penalty.

\subsection{Successive-term ratios}

For fixed $k$,
\begin{equation}\label{bsw:eq:ratio}
 \frac{B_{n+1,k}}{B_{n,k}}
 =\prod_{j=1}^k(kn+j)
 \left[1+\frac{(k-1)^2}{2kn^2}
 +\frac{(k-1)^2(k-2)}{6k^2n^3}
 +O_k(n^{-4})\right].
\end{equation}
Indeed, the factorial ratio is exact and subtraction of
$\log p_{n,k}$ from $\log p_{n+1,k}$ gives the bracket.  Such a ratio is
useful for checking recurrences attached to individual OEIS rows.

\subsection{A Lambert-\texorpdfstring{$W$}{W} inverse transseries}

Suppose a large value $Y$ is known and one wants to estimate the index $n$
for which $B_{n,k}\approx Y$.  Set
\begin{equation}\label{bsw:eq:Linv}
 L=\log Y+\lambda,
 \qquad
 N_0=\frac{L}{W_0(L/\ee)}.
\end{equation}
Then $N_0(\log N_0-1)=L$.  One Newton correction from the remaining terms in
\eqref{bsw:eq:logB} gives
\begin{equation}\label{bsw:eq:inverse}
 \boxed{\displaystyle
 N=N_0-
 \frac{\tfrac12\log(2\pi N_0)+A_1/N_0+A_2/N_0^2+A_3/N_0^3}
      {\log N_0}
 +O_k(N_0^{-1}),
 \qquad n=\frac Nk.}
\end{equation}
This is an inverse transseries in the same spirit as the inversion examples in
the ProveIt repository.  Further Newton steps and further terms of
Theorem~\ref{bsw:thm:all-orders} extend it mechanically.
[Added 1 October 2026, batch 73O2: the inversion examples meant are those of
the canonical volume \emph{Transseries and Inversion}~\cite{tai}.  The forward
map $\log B_{n,k}$ of \eqref{bsw:eq:logB} is the volume's gamma carrier
(theorem ``All-orders expansion of $\Gamma_+^{-1}$'', label
\texttt{p6:thm:gamma}) perturbed by the power series in $1/N$ of
\eqref{bsw:eq:logp}, the situation of its theorem ``Perturbed inversion''
(\texttt{p0:thm:perturbed-inversion}); the volume's chapter ``The
subfactorial'' (\texttt{p8:sec:top}) treats the derangement numbers the same
way.  Normalizations differ: here $L=\log Y+\lambda$ and $N_0=L/W_0(L/\ee)$,
with $\tfrac12\log(2\pi N_0)$ left in the correction; the volume uses
$R=\log y-\tfrac12\log(2\pi)$.  Formula \eqref{bsw:eq:inverse} is a
first-order instance of that apparatus, and no novelty is claimed for it.
Deciding the integer index from an approximate inverse, as in the deck
example below and in research question~11 of
Section~\ref{bsw:sec:questions}, is the subject of the volume's staircase
theorem (\texttt{p0:thm:staircase}).]

For the exact standard-deck count in \eqref{bsw:eq:deck-prob}, with $k=4$, the
leading Lambert estimate gives $n\approx13.17720934$.  The correction
\eqref{bsw:eq:inverse} gives
\[
 n\approx12.99979618,
\]
which correctly identifies the integer index $13$.

\section{Relation to rook polynomials and recent multipartite work}

The coefficients $a_{k,m}=m!\binom{k}{m}\binom{k-1}{m}$ are the familiar
rook numbers of a rectangular board of dimensions $k$ by $k-1$.  This
explains why rook-polynomial methods naturally occur in exact deck-shuffle
enumeration.  Enouen's report on the standard deck~\cite{enouen} develops such a direct
rook-polynomial approach.  The reciprocal polynomial $C_k(z)=R_k(-z)$ packages
all ranks into the compact power $C_k(z)^n$.

Balanced Smirnov words have also been connected explicitly with Hamiltonian
paths and cycles in complete multipartite graphs.  Recent work by
Mehiri~\cite{mehiri} gives
bijections and inclusion--exclusion formulas, including nonuniform part sizes.
The present asymptotic regime keeps the part size $k$ fixed and lets the number
of parts $n$ grow.  In this regime, Theorem~\ref{bsw:thm:factorial-moments} says:
\begin{quote}
Among all ordered vertex lists of $K_{k,k,\ldots,k}$, the fraction that form
Hamiltonian paths tends to $\ee^{-(k-1)}$, and the number of forbidden
within-part consecutive steps tends to $\Pois(k-1)$.
\end{quote}
This probabilistic statement complements exact Hamiltonian enumerations and
suggests analogous cyclic and nonuniform limit laws.

\section{Reproducibility and a formalization route}

The accompanying program \texttt{code/verify.py} performs the following tasks
using exact Python integers for enumeration and high-precision decimal
arithmetic for probabilities:
\begin{enumerate}[leftmargin=2em]
\item constructs $C_k(z)=\sum_m(-1)^ma_{k,m}z^m$;
\item raises it to the $n$th power by exact polynomial convolution;
\item evaluates \eqref{bsw:eq:grouped-sum};
\item checks the initial terms of A007060, A193624, and A330266;
\item verifies the exact 52-card value;
\item writes the numerical convergence table used above.
\end{enumerate}
The optional SymPy script \texttt{code/derive\_symbolic.py} reconstructs the
first four factorial cumulants and checks the displayed $n^{-3}$ expansion.
The article does not rely on a fitted numerical constant.

A Lean formalization can be separated into finite and analytic layers.
The finite core is especially approachable:
\begin{enumerate}[leftmargin=2em]
\item define rank words, runs, and the statistic $X$;
\item prove the run-compression bijection and finite coefficient identity
\eqref{bsw:eq:coeffA};
\item prove the polynomial identity \eqref{bsw:eq:Pexplicit};
\item reduce the integral notation to the factorial moment identity
$\int_0^\infty t^m\ee^{-t}dt=m!$;
\item prove the exact finite sum \eqref{bsw:eq:grouped-sum};
\item establish the leading factorial-moment estimate using polynomial degree
bounds.
\end{enumerate}
The final method-of-moments step can either use an existing probability
convergence library or be replaced by direct convergence of each fixed mass
through the probability generating function.  The exact count identities can
be formalized independently of the asymptotic layer.

\section{Proposed OEIS updates}

Subject to independent review and attribution, the following concise comments
would strengthen the relevant entries.

\paragraph{General fixed-$k$ family.}
If $B_{n,k}$ counts permutations of $kn$ labeled cards, $k$ in each of $n$
ranks, with no equal adjacent ranks, then for each fixed $k\ge2$,
\[
 \frac{B_{n,k}}{(kn)!}=\ee^{-(k-1)}\left(1-\frac{(k-1)^2}{2kn}
 +O_k(n^{-2})\right).
\]
More precisely, the number of equal-rank adjacencies converges to
$\Pois(k-1)$, and \eqref{bsw:eq:pfull} gives three correction terms.

\paragraph{A330266.}
Add \eqref{bsw:eq:k4}, together with the exact variance
$3-9/(4n-1)$ for the adjacency statistic.  The integral on the entry is the
$k=4$ instance of
\[
 B_{n,k}=\int_0^\infty
 \{(-1)^kk!L_k^{(-1)}(t)\}^n\ee^{-t}\,dt.
\]

\paragraph{A007060 and A193624.}
Add \eqref{bsw:eq:k2} and \eqref{bsw:eq:k3}, respectively.  Their multiset partners
A114938 and A193638 have the same normalized probabilities because of
\eqref{bsw:eq:labeled-multiset}.

A plain-text version is included as \texttt{oeis\_proposed\_updates.txt}; it
is explicitly marked as unsubmitted.
[Added 1 October 2026, batch 73O2: shipped as
\texttt{data/oeis\_proposed\_updates.txt}, unchanged.  Nothing has been
submitted to the OEIS.  Its correction terms rest on the sketched step of
Remark~\ref{bsw:rem:tail-sketch}.]

\section{Further research questions and topics}
\label{bsw:sec:questions}

\begin{enumerate}[leftmargin=2em]
\item \textbf{Uniformly growing multiplicity.}
Determine the largest range $k=k(n)$ in which a Poisson approximation remains
valid, and locate transitions to compound-Poisson or normal behavior.  The
fixed-$k$ proof is not uniform when the degree of $R_k$ grows.

\item \textbf{Explicit total-variation bounds.}
Use Stein--Chen coupling or the gamma representation to obtain a sharp bound
of the form
$d_{\mathrm{TV}}(X_{n,k},\Pois(k-1))\le C_k/n$, preferably with the optimal
constant and a computable second-order signed correction.

\item \textbf{Cyclic words and Hamiltonian cycles.}
Also forbid equality between the last and first ranks.  Find the exact
analogue of \eqref{bsw:eq:log-pgf-exp}, taking proper account of oriented versus
unoriented cycles and rotational symmetries.

\item \textbf{Unequal multiplicities.}
For ranks of sizes $k_1,\ldots,k_n$ and total $N$, identify natural sparse
conditions under which the adjacency count is asymptotically Poisson with
mean
$\sum_i k_i(k_i-1)/N$.  Determine the correction terms in symmetric sums of
the $k_i$.

\item \textbf{Larger exclusion distance.}
Forbid two equal ranks within distance $d$ rather than only distance one.
Cluster expansions should produce a Poisson or compound-Poisson law when
$k,d$ are fixed, but overlapping short patterns require new local cumulants.

\item \textbf{Moderate and large deviations.}
Analyze $\PP(X_{n,k}\ge x)$ outside the fixed-$x$ regime.  The exact gamma
expectation may lead to a saddle-point approximation uniformly for
$x$ growing with $n$.

\item \textbf{All-order closed forms.}
The operator rule \eqref{bsw:eq:operator-expansion} is constructive.  Find a
canonical Bell-polynomial expression for the coefficient of $n^{-r}$ in
$\log\Phi_{n,k}(v)$, and study sign or factorization patterns in $k$ and $v$.

\item \textbf{Laguerre-zero geometry.}
Clarify how the real zeros of $L_k^{(-1)}$ govern recurrences, cancellations,
and possible saddle points of the exact integral.  The integral involves a
sign-changing polynomial, so a naive positive Laplace principle is
insufficient.

\item \textbf{Differential equations and recurrences.}
For each fixed $k$, derive minimal-order polynomial recurrences for
$B_{n,k}$ and compare their asymptotic solutions with \eqref{bsw:eq:pfull}.  Seek
bounds on recurrence order as a function of $k$.

\item \textbf{New OEIS rows and arrays.}
Create a two-dimensional array for $S_{n,k}$ or $B_{n,k}$, with rows indexed
by $k$, including $k=5,6$ and exact kernels $P_k$.  Cross-reference the
labeled, multiset, rook-polynomial, and multipartite-graph interpretations.

\item \textbf{Rigorous inverse error bounds.}
Turn \eqref{bsw:eq:inverse} into an explicit bracketing theorem for the least
$n$ with $B_{n,k}\ge Y$, including monotonicity thresholds and certified
rounding.

\item \textbf{Machine-checked asymptotics.}
Formalize the finite identities in Lean, then isolate a reusable theorem that
turns fixed factorial-moment convergence plus tightness into Poisson
convergence.  This would support many related OEIS limit laws.
\end{enumerate}

\section{Conclusion}

The constant $\ee^{-(k-1)}$ is not an isolated feature of three OEIS rows.  It
is the zero mass of a limiting Poisson law for the complete adjacency
statistic.  The run transform converts the combinatorics into a single
Laguerre polynomial; reciprocal coefficients give exact inclusion--exclusion;
and a gamma negative-moment representation produces systematic correction
terms.  The resulting framework simultaneously proves the fixed-$k$
conjecture recorded on A330266, sharpens the known $k=2,3,4$ asymptotics, and
opens a tractable family of cyclic, nonuniform, and growing-multiplicity
problems.

\appendix
\section{Cumulant algebra}\label{bsw:app:cumulant-algebra}

For completeness, write
\[
 R_k(z)=1+a_1z+a_2z^2+a_3z^3+a_4z^4+O(z^5),
\]
where
\begin{align*}
 a_1&=k(k-1),\\
 a_2&=\frac{k(k-2)(k-1)^2}{2},\\
 a_3&=\frac{k(k-3)(k-2)^2(k-1)^2}{6},\\
 a_4&=\frac{k(k-4)(k-3)^2(k-2)^2(k-1)^2}{24}.
\end{align*}
The coefficient of $v^r$ in \eqref{bsw:eq:Phi-exact} is
$\coef{z^r}R_k(z)^n/\fall Nr$.  Thus, for example,
\begin{align*}
 \EE\fall X2
 &=\frac{2na_2+n(n-1)a_1^2}{\fall N2},\\
 \EE\fall X3
 &=\frac{6na_3+6n(n-1)a_2a_1
       +n(n-1)(n-2)a_1^3}{\fall N3}.
\end{align*}
The fourth factorial moment is obtained from the five partitions
$4$, $3+1$, $2+2$, $2+1+1$, and $1+1+1+1$.  Applying the standard
moment-to-cumulant relations and simplifying gives
\eqref{bsw:eq:kappa1}--\eqref{bsw:eq:kappa4}.  The accompanying SymPy script performs
exactly this finite calculation and verifies that the difference from
\eqref{bsw:eq:log-pgf-exp} is $O(n^{-4})$.

\section{Exact initial terms for \texorpdfstring{$k=5$}{k=5}}

The same formula supplies a useful prospective OEIS row.  For $n=0,1,\ldots,7$,
\begin{align*}
 B_{n,5}={}&1,\ 0,\ 28800,\ 12420864000,\ 27917203599360000,\\
 &197726965332480000000000,\\
 &3614614673967242167910400000000,\\
 &147444323126376251733903684403200000000.
\end{align*}
The kernel is
$P_5(t)=t^5-20t^4+120t^3-240t^2+120t$, and
$B_{n,5}=\int_0^\infty P_5(t)^n\ee^{-t}\,dt$.
These terms should be independently cross-checked before creating or editing
an OEIS entry.

\section{Provenance and editorial notes}\label{bsw:app:provenance}

[Added 1 October 2026, batch 73O2.]

\textbf{Source.}  The report was built from one manuscript, manuscript~57 of
batch~73 of ProveIt's incoming-reports intake: the archive
\texttt{ProveIt\_Balanced\_Smirnov\_OEIS.zip} (main file
\texttt{article.tex}, a 17-page PDF).  It arrived in commit
\texttt{aa43cc555} and was placed in commit \texttt{6e193dd4f}, which deleted
the archive from \texttt{docs/incoming}; it survives in \texttt{aa43cc555}.
The manuscript pins no repository revision and cites no repository file.  It
has no author or preparer line: its author field is the subtitle ``A research
note on OEIS A007060, A193624, A330266, and related sequences'', and the
delivered PDF's author metadata is empty.

\textbf{What the write changed.}
\begin{itemize}[leftmargin=2em]
\item Every label gained the prefix \texttt{bsw:}; none was otherwise renamed
 or removed.  The manuscript had 66 labels.
\item Citation commands were added: the manuscript listed twelve
 bibliography entries but cited none of them.  They now stand at the first
 mentions of the OEIS entries, the Smirnov-word generating function, the
 Laguerre expansion, and the works of Eriksson--Martin, Enouen and Mehiri.
\item The authors and dates of the two OEIS contributions are named:
 S.~Kirgizov's fixed-$k$ conjecture (29 September 2023) and D.~Radcliffe's
 deduction of the case $k=4$ (9 September 2025).
\item Added: the editorial note after the scope paragraph,
 Remark~\ref{bsw:rem:tail-sketch} on the sketched analytic step, a note in
 the inverse section citing the transseries volume~\cite{tai}, a note on the
 shipped file of proposed OEIS comments, one
 bibliography entry, the label \texttt{bsw:sec:questions}, and this section.
\item Nothing was removed.  Every non-claim of the manuscript is kept: no
 complete literature search or priority claim, the numerical residuals ``not
 used as proof'', the cross-check required for the $k=5$ terms, and the
 unsubmitted status of the proposed OEIS comments.
\end{itemize}

\textbf{Relation to the companion report.}  The path-forest report
\texttt{a189281-path-forest-expansions} of this collection (batch~73,
manuscript~60) follows the same chain of method (exact marked-subset
generating function, Poisson limit, all-orders $1/n$ expansion of a
factorial probability generating function, Lambert-$W$ inversion) for
permutations avoiding a fixed positional and value offset, with limit
Poisson$(\theta)$, $\theta\in\{1,2\}$.  Neither theorem specializes to the
other.  In the present model the $n$ ranks give $n$ cliques $K_k$ as the
``value'' relation, whose number grows with $n$, so the hypotheses of the
path-forest report (a fixed number of long paths) do not hold.  Its uniform
factorial-moment bound is the device named in
Remark~\ref{bsw:rem:tail-sketch}.

\textbf{Checks by the intake.}  The intake reran both shipped programs on a
copy (the CSV is byte-identical; the deck output equal apart from line
endings; the symbolic script confirms $\kappa_1^{(F)},\ldots,\kappa_4^{(F)}$
and the $n^{-3}$ expansion), compared the terms of A007060, A193624 and
A330266 with an independent dynamic program over rank words, confirmed the
$k=5$ terms (for $n\le4$ by that program, for $n\le7$ by the formula) and the
52-card value, checked the variances $4/5$, $6/5$, $6/7$ at
$(n,k)=(3,2),(2,3),(4,2)$ by full enumeration, and rederived by hand the
local correction, the ratio bracket and $A_1,A_2,A_3$.  These checks do not
replace an independent review of the proofs.

\begin{thebibliography}{99}

\bibitem{oeis330266}
OEIS Foundation Inc.,
\emph{A330266: Number of ways to shuffle a deck of $4n$ cards, with four
cards in each of $n$ ranks, so that adjacent cards have different ranks},
\url{https://oeis.org/A330266}, accessed 1 October 2026.

\bibitem{oeis007060}
OEIS Foundation Inc.,
\emph{A007060: Number of ways $n$ married couples can sit in a row without
any spouses next to each other},
\url{https://oeis.org/A007060}, accessed 1 October 2026.

\bibitem{oeis193624}
OEIS Foundation Inc.,
\emph{A193624: Number of ways $n$ triples can sit in a row without any
siblings next to each other},
\url{https://oeis.org/A193624}, accessed 1 October 2026.

\bibitem{oeis114938}
OEIS Foundation Inc.,
\emph{A114938: Permutations of $\{1,1,2,2,\ldots,n,n\}$ with no equal
consecutive terms},
\url{https://oeis.org/A114938}, accessed 1 October 2026.

\bibitem{oeis193638}
OEIS Foundation Inc.,
\emph{A193638: Permutations of the balanced multiplicity-three multiset with
no equal consecutive terms},
\url{https://oeis.org/A193638}, accessed 1 October 2026.

\bibitem{oeis321633}
OEIS Foundation Inc.,
\emph{A321633: Permutations of the balanced multiplicity-four multiset with
no equal consecutive terms},
\url{https://oeis.org/A321633}, accessed 1 October 2026.

\bibitem{eriksson-martin}
H.~Eriksson and A.~Martin,
\emph{Enumeration of Carlitz multipermutations},
arXiv:1702.04177, 2017.

\bibitem{enouen}
J.~Enouen,
\emph{What is the Perfect Shuffle?},
arXiv:1911.07426, 2019.

\bibitem{mehiri}
E.-M.~Mehiri,
\emph{Bijections Between Smirnov Words and Hamiltonian Cycles in Complete
Multipartite Graphs}, arXiv:2510.26597v2, 2026.

\bibitem{stanley}
R.~P.~Stanley,
\emph{Enumerative Combinatorics, Volume 1}, second edition,
Cambridge University Press, 2012.

\bibitem{flajolet-sedgewick}
P.~Flajolet and R.~Sedgewick,
\emph{Analytic Combinatorics}, Cambridge University Press, 2009.

\bibitem{dlmf}
NIST Digital Library of Mathematical Functions,
chapters on the gamma and Laguerre functions,
\url{https://dlmf.nist.gov/}, accessed 1 October 2026.

\bibitem{tai}\begingroup\raggedright
[Added in the batch-73O2 write.] V.~Reshetnikov and contributors,
\emph{Transseries and Inversion}, canonical volume of the ProveIt
series-and-transseries collection.  Its source is
\path{Analysis/Transseries/docs/series-and-transseries/}\allowbreak
\path{Transseries_And_Inversion/transseries_and_inversion.tex}.  Cited here:
Part ``Inversion: the apparatus for a rapidly growing function''
(labels \texttt{p0:thm:perturbed-inversion}, \texttt{p0:thm:staircase}),
chapter on the increasing inverse of Gamma (\texttt{p6:thm:gamma}) and chapter
``The subfactorial'' (\texttt{p8:sec:top}).  Repository state of placement
commit \texttt{6e193dd4f}.\par\endgroup

\end{thebibliography}

\end{document}
